\subsection{Channel Ablation Study}
\label{sec:abla-study}

To isolate and study the contribution of individual channels, we evaluate SS-JEPA pre-trained on different subsets of channels, with a UPerNet~\cite{DBLP:journals/corr/abs-1807-10221} segmentation decoder, in \cref{tab:channel_ablation}.

Utilizing purely optical channels (RGB + Grayscale) establishes a baseline performance of $0.812$ mIoU. This metric highlights a persistent limitation in planetary surface imaging, where atmospheric dust variations and uniform albedo profiles frequently obscure landform boundaries. In contrast, isolating structural geometry channels (DEM + Slope) improves performance to $0.835$ mIoU. This margin demonstrates that topographic markers, specifically the abrupt elevation drops found at source scarps and the adjacent depositional accumulation zones, provide highly discriminative features that remain entirely invariant to surface illumination or optical noise. Progressive cross-modal integration yields consistent performance gains, with the full 7-channel configuration achieving a peak score of 0.865 mIoU. While the thermal inertia modality exhibits a lower native spatial resolution, its integration with geometric features provides a $0.011$ mIoU performance boost. These results demonstrate that combining optical, topographic, and thermophysical data allows our framework to extract complementary physical structures, maximizing landslide segmen accuracy across heterogeneous Martian terrains.

\begin{table}
	\centering
	\begin{tabular}{cccccc}
		\toprule
		\textbf{RGB} & \textbf{Gray} & \textbf{DEM} & \textbf{Slope} & \textbf{Thermal} & \textbf{mIoU}  \\
		\midrule
		\checkmark   & \checkmark    &              &                &                  & 0.812          \\
		             &               & \checkmark   & \checkmark     &                  & 0.835          \\
		\checkmark   & \checkmark    & \checkmark   &                &                  & 0.841          \\
		\checkmark   & \checkmark    & \checkmark   & \checkmark     &                  & 0.854          \\
		             &               & \checkmark   & \checkmark     & \checkmark       & 0.846          \\
		\midrule
		\checkmark   & \checkmark    & \checkmark   & \checkmark     & \checkmark       & \textbf{0.865} \\
		\bottomrule
	\end{tabular}
	\vspace{0.5em}
	\caption{Channel Ablation Study of the SS-JEPA + UPerNet Segmenation Model Configuration}
	\label{tab:channel_ablation}
\end{table}


